\chapter{Export fidelity and the featuriser fallback}

The phase that expected a fight with a converter and got a different problem
instead.

\begin{defect}{The fallback was the design, not a retreat}
\symptom{The plan anticipated a fight with the converter for text vectorizers and
named a fallback: move the whole text pipeline into hand written code, export
only the linear layer, and take that quickly rather than spending a day coaxing a
converter.}
\diagnosis{The fight never started, because the feature set was never
convertible. Three of the five feature blocks are categorical one hots, option
shape booleans and structural buckets, and no standard transformer produces
those. Putting the whole featurisation in the graph would have meant writing
custom converters for three bespoke transformers \emph{and} trusting the
vectorizer converters, which is a larger version of exactly the risk the fallback
exists to avoid.}
\resolution{The hand written featuriser, shared by training and inference. It is
the only design in which one implementation serves both paths, which the feature
specification required from the beginning.}
\instead{Writing the custom converters. The second argument against arrived for
free and will matter longer: a vectorizer anywhere in the inference path puts a
full machine learning framework on the dependency list of a tool that installs
with one command. The route taken means an installed package needs an array
library and a runtime and nothing else.}
\end{defect}

\begin{defect}{A reimplementation that is a claim rather than an identity}
\symptom{A hand written character n-gram analyser is slower than a compiled one,
and its faithfulness to the reference analyser is an assertion.}
\diagnosis{So assert it, in a test, rather than in a comment.}
\resolution{A test asserts, over inputs including empty strings, words shorter
than the n-gram length, and Japanese text, that this project's enumeration is
character for character the reference implementation's.}
\instead{Spot checking a few strings. The detail an independent reimplementation
gets wrong is the short word rule: a word shorter than the n-gram length yields
its padded self once rather than once per length. Getting that wrong silently
triples the weight of every two letter token, and it would not show up in a spot
check of ordinary words.}
\end{defect}

\begin{defect}{An opset probe, and a warning that is recorded rather than
suppressed}
\symptom{The export is specified to use the newest opset the resolved runtime
accepts, determined by probe. The probe walks down from the newest the installed
format library defines and takes the first that both converts and opens in a
session. The converter warns at every attempt that the requested opset is above
the newest it has been tested against.}
\diagnosis{The warning is not noise. It is precisely the case where a successful
export proves nothing: the converter is saying it has no test coverage for the
combination it just produced.}
\resolution{The warning is recorded in the decision record rather than filtered,
and the thing that proves something is the parity gate. That is why the export
phase has a gate rather than a deliverable.}
\instead{Suppressing the warning and trusting the successful conversion. A
conversion that succeeds and produces different numbers is the failure mode the
whole gate exists for.}
\end{defect}

\begin{defect}{The graph cannot read its own weights}
\symptom{Naming the n-grams behind a prediction, which is what the evidence
requirement in this project's first law demands, could not be done from the
inference path.}
\diagnosis{The exported graph is one classifier node followed by a normaliser,
and the coefficients, the intercepts and the class order live in the node's
\emph{attributes} rather than in graph initialisers. The runtime does not hand a
caller a node's attributes. Reading them would have required the full model
format library at runtime to produce one line of a report.}
\resolution{A separate evidence table, holding each class's highest weighted
features, written by the same training run that produces the model.}
\instead{Adding the model format library as a runtime dependency, or dropping the
evidence line from findings. The first inflates every install for one line of
output; the second violates the law the tool is built around.

A derived file is a drift risk, so the parity suite reads the weights back out of
the graph and asserts that the table agrees with them. A stale evidence file now
fails the build rather than naming the wrong n-grams in a finding, which would
have satisfied the evidence law in form and violated it in substance.}
\end{defect}

\begin{defect}{The first training run was thrown away}
\symptom{The manifest of the first full training run recorded that the working
tree was dirty.}
\diagnosis{The reproducibility rules say a result produced from a dirty tree is
marked dirty and is not citable. The flag did exactly what it exists for.}
\resolution{The run was discarded and repeated from a clean tree.}
\instead{Citing it with a note. The flag needed one fix of its own to be honest:
the output directory is untracked on a first run, so counting it would have
marked every first training run dirty by construction and made the flag mean
nothing. It now excludes the output directory and only that.

This entry has a sequel. See the open item in
Chapter~\ref{ch:still-open} about the training manifest of the currently shipped
model, which is marked dirty for a different and less avoidable reason and was
not thrown away.}
\end{defect}
